\textbf{Задание №7} Подготовить отчёт по работе, который должен включать: краткое описание всех использованных системных вызовов и команд; текст программы; протоколы; скриншоты с демонстрацией поведения ОС Linux; выводы. \label{sec:7}

    Ниже перечислены все системные вызовы, библиотечные функции и структуры, задействованные в программе «Диспетчер с закрытым окном», с кратким описанием их назначения.

    \begin{enumerate}

        \item \textbf{Управление процессами}

        \begin{enumerate}
            \item \texttt{fork()} — создаёт новый процесс-потомок, копирующий адресное пространство родителя (механизм copy-on-write). Возвращает PID потомка в родителе, $0$ в потомке и $-1$ при ошибке. Используется для порождения $N$ исполнителей.
            \item \texttt{\_exit(status)} — немедленно завершает процесс, минуя обработчики \texttt{atexit} и очистку буферов stdio. Применяется в детях при получении \texttt{SIGTERM}.
            \item \texttt{getpid()} — возвращает PID текущего процесса (используется детьми для инициализации \texttt{srand}).
            \item \texttt{getppid()} — возвращает PID родительского процесса; в детях применяется как адрес получателя сигналов \texttt{SIGUSR1}/\texttt{SIGUSR2}.
            \item \texttt{waitpid(pid, *status, options)} — ожидает завершения конкретного потомка. В коде вызывается с \texttt{options = 0} для синхронного ожидания после рассылки \texttt{SIGTERM}.
        \end{enumerate}

        \item \textbf{Работа с сигналами}

        \begin{enumerate}
            \item \texttt{signal(signum, handler)} — устанавливает обработчик сигнала. В родителе применяется для \texttt{SIGUSR1}, \texttt{SIGUSR2}, \texttt{SIGINT} и \texttt{SIGALRM}; в детях — для \texttt{SIGTERM}. Также используется \texttt{SIG\_IGN} для игнорирования \texttt{SIGINT}, \texttt{SIGUSR1} и \texttt{SIGUSR2} в дочерних процессах.

            \item \texttt{kill(pid, sig)} — посылает сигнал процессу с указанным PID. Дети шлют родителю \texttt{SIGUSR1}/\texttt{SIGUSR2}; родитель посылает детям \texttt{SIGTERM} при завершении.

            \item \texttt{sigemptyset(set)} — инициализирует набор сигналов пустым множеством.

            \item \texttt{sigaddset(set, signum)} — добавляет сигнал в набор. В коде добавляются \texttt{SIGUSR1} и \texttt{SIGUSR2} в маску блокировки.

            \item \texttt{sigprocmask(how, set, oldset)} — изменяет маску блокируемых сигналов:
            \begin{enumerate}
                \item \texttt{SIG\_BLOCK} — добавляет сигналы в блокировку (начало закрытого окна);
                \item \texttt{SIG\_SETMASK} — восстанавливает прежнюю маску (конец окна).
            \end{enumerate}

            \item \texttt{sigpending(set)} — заполняет набор сигналами, находящимися в состоянии ожидания доставки (в pending). В закрытом окне показывает наличие заблокированных сигналов.

            \item \texttt{sigismember(set, signum)} — проверяет, входит ли сигнал в набор. Используется для вывода факта наличия \texttt{SIGUSR1}/\texttt{SIGUSR2} в pending (значение $0$ или $1$, а не количество).

            \item \texttt{sigsuspend(mask)} — атомарно заменяет маску сигналов процесса и приостанавливает его до получения любого сигнала, после чего восстанавливает исходную маску. Это единственный способ ожидания сигнала, исключающий гонку между проверкой условия и переходом в сон. В коде вызывается в открытом окне с пустой маской \texttt{empty}.

            \item \texttt{alarm(seconds)} — устанавливает таймер, по истечении которого процессу будет послан сигнал \texttt{SIGALRM}. Используется для гарантированного пробуждения родителя из \texttt{sigsuspend} раз в секунду, чтобы обновить проверку времени.
        \end{enumerate}

        \item \textbf{Управление памятью и разделяемое состояние}

        \begin{enumerate}
            \item \texttt{mmap(...)} — отображает разделяемую анонимную память (\texttt{MAP\_SHARED | MAP\_ANONYMOUS}) для счётчиков отправленных сигналов, доступных одновременно родителю и всем потомкам.
            \item \texttt{munmap(...)} — освобождает ранее отображённую область памяти перед завершением программы.
            \item \texttt{memset(...)} — обнуляет содержимое выделенной разделяемой памяти при инициализации.
        \end{enumerate}

        \item \textbf{Время и задержки}

        \begin{enumerate}
            \item \texttt{sleep(sec)} — приостанавливает процесс на указанное число секунд (закрытое окно).
            \item \texttt{usleep(usec)} — микросекундная задержка: интервал между сигналами у детей (случайный, $100$--$400$ мс) и пауза после разблокировки для доставки сигналов.
            \item \texttt{alarm(1)} — устанавливает сигнал \texttt{SIGALRM} через одну секунду; применяется для периодического пробуждения родителя из \texttt{sigsuspend} и проверки длительности открытого окна.
            \item \texttt{time(NULL)} — возвращает текущее время в секундах; используется для контроля длительности открытого окна ($\texttt{OPEN\_SEC}$).
        \end{enumerate}

        \item \textbf{Генерация случайных чисел}

        \begin{enumerate}
            \item \texttt{srand(seed)} — инициализирует генератор псевдослучайных чисел; в детях используется \texttt{getpid()}, чтобы у каждого процесса была своя последовательность.
            \item \texttt{rand()} — возвращает случайное число; применяется для выбора типа сигнала (\texttt{rand() \& 1}) и случайной задержки ($100 + \texttt{rand()\%300}$ мс).
        \end{enumerate}

        \item \textbf{Ввод-вывод и стандартные функции}

        \begin{enumerate}
            \item \texttt{setvbuf(stdout, NULL, \_IONBF, 0)} — отключает буферизацию стандартного вывода, чтобы диагностические сообщения появлялись в консоли сразу.
            \item \texttt{printf(...)} — форматированный вывод статистики (начало и конец открытого окна, состояние pending, $\Delta$ обработанных сигналов, итоговая статистика).
        \end{enumerate}

        \item \textbf{Пользовательские типы и макросы}

        \begin{enumerate}
            \item \texttt{volatile sig\_atomic\_t} — тип, безопасно изменяемый в обработчиках сигналов. Используется для счётчиков \texttt{sent1}, \texttt{sent2} (в разделяемой памяти), а также \texttt{recv1}, \texttt{recv2}, \texttt{stop}.
            \item \texttt{pid\_t} — тип идентификатора процесса.
            \item \texttt{sigset\_t} — тип набора сигналов, применяемый в \texttt{sigemptyset}, \texttt{sigaddset}, \texttt{sigprocmask}, \texttt{sigpending}, \texttt{sigismember}, \texttt{sigsuspend}.
            \item \texttt{sigismember} возвращает ненулевое значение при наличии сигнала в наборе и нуль в противном случае.
        \end{enumerate}

        \item \textbf{Общая схема работы}

        \begin{enumerate}
            \item Родитель порождает $N = 3$ дочерних процесса через \texttt{fork()}. Дети наследуют код, но получают собственные PID и адресные пространства.

            \item Дети настраивают обработчики: \texttt{SIGINT}, \texttt{SIGUSR1}, \texttt{SIGUSR2} игнорируются, а \texttt{SIGTERM} завершает процесс через \texttt{\_exit(0)}. Затем каждый ребёнок в бесконечном цикле со случайной задержкой $100$--$400$ мс посылает родителю случайный сигнал (\texttt{SIGUSR1} или \texttt{SIGUSR2}), инкрементируя соответствующий счётчик в разделяемой памяти.

            \item Родитель выполняет $3$ цикла, каждый из которых состоит из двух фаз:
            \begin{enumerate}
                \item \emph{Открытое окно} ($\texttt{OPEN\_SEC} = 2$ с): сигналы не заблокированы, обработчики срабатывают сразу. Ожидание реализовано через \texttt{sigsuspend(\&empty)} с периодическим \texttt{alarm(1)}, что исключает гонку между проверкой условия и переходом в сон.
                \item \emph{Закрытое окно} ($\texttt{CLOSE\_SEC} = 3$ с): вызывается \texttt{sigprocmask(SIG\_BLOCK, ...)}, и \texttt{SIGUSR1}/\texttt{SIGUSR2} копятся в pending. По окончании \texttt{sleep()} вызывается \texttt{sigpending()}, и с помощью \texttt{sigismember()} выводится факт наличия сигналов в pending. Затем маска восстанавливается через \texttt{SIG\_SETMASK}, и доставляются отложенные сигналы.
            \end{enumerate}

            \item Сравнение \texttt{sent} (из разделяемой памяти) и \texttt{recv} (счётчик обработчиков) демонстрирует, что обычные сигналы не ставятся в очередь: за время блокировки в pending всегда находится только один бит на каждый тип, поэтому после разблокировки доставляется не более одного-двух экземпляров сигнала. Итоговые потери достигают $70$--$80\%$.

            \item По \texttt{SIGINT} родитель устанавливает флаг \texttt{stop = 1}, игнорирует дальнейшие \texttt{SIGINT}, рассылает детям \texttt{SIGTERM} и собирает их через \texttt{waitpid()}. Это обеспечивает корректное завершение без зомби-процессов и утечек ресурсов.
        \end{enumerate}

        \newpage
        
        \item \textbf{Экспериментальные результаты}

        \begin{enumerate}
            \item При $\texttt{N} = 3$, $\texttt{OPEN} = 2$ с, $\texttt{CLOSE} = 3$ с за три цикла было отправлено $99 + 78 = 177$ сигналов, а обработано $19 + 19 = 38$, то есть потеряно около $78\%$.

            \item В каждом закрытом окне \texttt{sigpending()} показывал \texttt{USR1=1}, \texttt{USR2=1} независимо от количества накопленных сигналов — это ключевое доказательство того, что pending хранит бит, а не счётчик.

            \item Приращение $\Delta$ после разблокировки всегда лежало в диапазоне $1$--$2$: один экземпляр из pending плюс (возможно) один свежий сигнал, пришедший сразу после сброса бита.

            \item Эти данные полностью согласуются с моделью POSIX: для обычных сигналов (не real-time) повторные доставки одного типа сливаются в один pending-бит, тогда как сигналы реального времени (\texttt{SIGRTMIN}\dots\texttt{SIGRTMAX}), передаваемые через \texttt{sigqueue()}, ставятся в очередь и не теряются.
        \end{enumerate}
    \end{enumerate}